\chapter{测试与性能分析}
\label{chap:testing}

\section{测试范围与数据集}

测试围绕三个功能阶段展开：
任务一检查五类文件路由、PDF 解析适配和原有格式兼容性；
任务二检查记录解析、级联抽取、可靠性分级和图谱入库；
任务三使用八个在线业务场景验证语义规划、只读查询、结果组织和统计图生成，
并通过 MCP 传输层核对统一分析入口和兼容入口的结果契约。
任务二当前提供 32 项独立自动化回归用例；任务一、任务三按各自评测脚本另行执行；
业务查询通过公开工作台接口连续请求三次并记录返回行数与响应时间。

\section{智能体工具调用核验}

智能体行为核验关注三个可观察结果：
工具选择是否与用户目标一致，
调用是否使用了真实的数据或查询参数，
返回结果是否包含回答所需的结构化事实。
分别执行统一医学分析查询和数据来源查询，
前者选择 \path{ask_medical_analytics}，
后者选择 \path{get_medical_data_sources}。
两次调用均返回对应的结构化数据，
最终回答仅采用工具返回的数量和医学事实。

\begin{table}[H]
\centering
\caption{代表性智能体工具路由核验}
\label{tab:agent-route-check}
\begin{tabularx}{\textwidth}{L{3.3cm}L{4.0cm}Y}
\toprule
\textbf{用户目标} & \textbf{实际选择的工具} & \textbf{可核对结果} \\
\midrule
查询肺不张症状 & \path{ask_medical_analytics} &
返回确定性查询计划、20 行证据、SQL 和分析追溯信息 \\
查询当前登记来源 & \path{get_medical_data_sources} &
返回来源总数、来源列表和记录规模 \\
\bottomrule
\end{tabularx}
\end{table}

工具路由与后端处理保持分层：
智能体只选择任务级入口，
工具服务负责参数校验和状态整理；
任务三的查询、图表和导出统一进入分析服务，
后续的算子执行、图谱事务或只读查询由确定性服务完成。
因此，路由核验与任务一、任务二和任务三的功能测试分别检查不同层次。

运行态还通过 MCP 的 Streamable HTTP 接口列出 19 个工具，
分别调用 \path{ask\_medical\_analytics} 和兼容的
\path{execute\_nl2sql}。
同一问题两次均返回成功、确定性 NL2SQL 规划和 20 行结果，
两次生成的 SQL 相同，并同时包含来源范围和执行追溯字段。
工作台 \path{/api/query} 也返回相同的分析状态和 20 行证据。
Nexent Agent 对话已实际选择 \path{ask\_medical\_analytics}，
但当前事件流只返回最终文本，没有返回可验证的结构化分析对象；
网关因此使用同一分析服务回退并标记降级。
这证明了服务链路统一，不把该次 Agent 对话描述为结构化事件闭环通过。

\section{任务一输入路由测试}

\repoPath{tests/task1/test_pdf_support.py} 包含 9 项自动化用例，
覆盖 PDF 接入新增逻辑及 TXT、CSV、JSON、JSONL 原有路由。
测试使用模拟的 MinerU 响应，
检查本地接口转换、格式路由与产物转换逻辑。
全部用例执行通过，分组结果见表~\ref{tab:unit-test-results}。

\begin{table}[H]
\centering
\caption{任务一格式与 PDF 路由测试结果}
\label{tab:unit-test-results}
\begin{tabularx}{\textwidth}{L{3.1cm}Y C{2.0cm}}
\toprule
\textbf{测试分组} & \textbf{覆盖内容} & \textbf{结果} \\
\midrule
格式路由与兼容性 &
PDF 分类、二进制预览保护、原有四类格式分组、
结构化处理链保持、PDF 转文本后复用文本链 &
5/5 通过 \\
\addlinespace
MinerU 接口适配 &
签名上传与任务轮询、结果下载、Markdown 正文转换、
TXT 产物与解析信息生成 &
3/3 通过 \\
\addlinespace
PDF 服务可用性路由 &
远程解析可用性识别与处理链选择 &
1/1 通过 \\
\midrule
合计 & 9 项测试 & 9/9 通过 \\
\bottomrule
\end{tabularx}
\end{table}

这组用例覆盖的结果为：
新增 PDF 分支保持了四类原有格式的分组和结构化处理链，
解析产物能够转为 TXT 并继续进入既有文本清洗流程，
上传、轮询和下载状态也能够转换为任务一统一结果。

\section{任务二抽取回归测试}

任务二的 32 项用例覆盖级联抽取和工具层的异常保护。
其中 8 项检查模型服务降级、空结果和未知可靠性等级的安全返回，
18 项检查长实体匹配、分段疾病归属、统一实体门禁和缺口候选，
6 项检查离线优先、批处理范围、复核决定和未复核候选的入库状态。
测试使用固定文本和模拟模型返回，
不依赖外部模型服务；本组 32 项用例在同一代码仓环境中全部通过。

\begin{table}[H]
\centering
\small
\caption{任务二自动化测试结果}
\label{tab:task2-test-results}
\begin{tabularx}{\textwidth}{L{3.0cm}Y C{2.0cm}}
\toprule
\textbf{测试分组} & \textbf{覆盖内容} & \textbf{结果} \\
\midrule
工具防护与降级 &
模型服务异常、空结果、未知可靠性等级和混合方式降级 & 8/8 通过 \\
\addlinespace
级联语义回归 &
长实体优先、分段疾病归属、统一实体门禁与原文范围保持 & 18/18 通过 \\
\addlinespace
级联流程与入库 &
离线优先、复核队列、批处理边界和未复核候选处理 & 6/6 通过 \\
\midrule
合计 & 32 项测试 & 32/32 通过 \\
\bottomrule
\end{tabularx}
\end{table}

\section{任务二固定留出集评测}

功能回归用于检查处理链是否按预期运行，质量评测则采用固定留出集计算实体和关系抽取指标。
本次评测使用 CMeEE 的 2\,500 条实体样本和 CMeIE 的 1\,793 条关系样本，
采用严格匹配统计精确率、召回率和 F1，并同步记录单条记录的中位耗时与 P95 耗时。
评测后端为 \texttt{offline}，结果对应词典、规则和版本配置均已固定。

\begin{table}[H]
\centering
\small
\caption{任务二固定留出集离线评测结果}
\label{tab:task2-holdout}
\begin{tabularx}{\textwidth}{L{2.1cm}L{2.3cm}C{1.1cm}C{1.2cm}C{1.2cm}C{1.2cm}Y}
\toprule
\textbf{评测对象} & \textbf{数据集} & \textbf{样本数} & \textbf{精确率} & \textbf{召回率} & \textbf{F1} & \textbf{耗时/ms（P50/P95）} \\
\midrule
实体识别 & CMeEE 固定留出集 & 2\,500 & 53.94\% & 28.97\% & 37.70\% & 2.26 / 33.84 \\
关系抽取 & CMeIE 固定留出集 & 1\,793 & 74.34\% & 7.92\% & 14.31\% & 7.84 / 48.21 \\
\bottomrule
\end{tabularx}
\end{table}

实体识别结果反映词典匹配对高频术语的覆盖，关系抽取结果反映当前离线规则在严格关系匹配下的有效输出。
两项指标共同作为任务二离线基线，后续模型增强评测沿用相同样本划分和匹配规则，保证结果可比较。

固定留出集用于正式报告；开发过程中另取 CMeEE、CMeIE 各 100 条做同批工程回归，
用于快速判断级联修改是否同时改善离线与混合链。该小样本不替代固定留出集成绩。
统一低可靠实体与模型补漏实体的门禁后，实体识别的混合方式达到
63.93\% 精确率、58.42\% 召回率和 61.05\% F1，
均高于同批离线方式的 60.80\%、39.19\% 和 47.66\%。
关系抽取增加分段疾病标题继承后，离线方式为
78.43\%/12.99\%/22.28\%，混合方式为
79.31\%/14.94\%/25.14\%；模型补充的关系保持严格证据门禁，
本轮 100 条记录未出现模型调用错误，也未写入知识图谱或刷新分析库。

\section{知识库与分析库核验}

公开工作台的数据概览返回 14\,408 个疾病条目、79\,605 个实体、
445\,820 条三元组、445\,210 条疾病事实、3\,836 条质量记录和 6 个来源。
基础结构化知识、CMeIE 关系与 DataMate 增量事实之和为 445\,820，
与图谱三元组总数一致。
分析数据库按疾病事实口径生成 445\,210 条记录。

\begin{table}[H]
\centering
\caption{知识库与分析库检查结果}
\label{tab:runtime-check}
\begin{tabularx}{0.86\textwidth}{Y C{3.0cm} C{2.2cm}}
\toprule
\textbf{检查项目} & \textbf{数量或状态} & \textbf{结果} \\
\midrule
图谱实体 & 79\,605 & 可查询 \\
图谱三元组 & 445\,820 & 可查询 \\
来源分项合计 & 445\,820 & 与总数一致 \\
疾病事实 & 445\,210 & 可查询 \\
分析成果导出 & HTML、CSV、JSON & 可用 \\
\bottomrule
\end{tabularx}
\end{table}

\section{在线查询验证}

业务查询覆盖疾病的症状、检查、治疗、科室和药物五类属性，
以及科室分布、实体类型分布和关系类型分布三类统计问题。
每个问题在同一公开工作台接口上连续请求三次，
表~\ref{tab:online-query-results} 给出页面实际展示行数、
三次请求的中位响应时间和波动范围。
响应时间包含 HTTPS 传输、语义规划、SQLite 查询和结果组织；
表中“范围”为三次请求的最小值至最大值，
本表用于描述该服务实例的响应情况，解释范围限定为当前运行配置。

\begin{table}[H]
\centering
\small
\caption{在线业务查询与响应时间}
\label{tab:online-query-results}
\begin{tabularx}{\textwidth}{L{3.0cm}Y C{1.5cm} C{1.8cm} C{2.1cm}}
\toprule
\textbf{场景} & \textbf{测试问题} & \textbf{行数} &
\textbf{中位/ms} & \textbf{范围/ms} \\
\midrule
疾病症状 & 肺炎有哪些症状？ & 40 & 719 & 637--1\,028 \\
疾病检查 & 肺炎应该做哪些检查？ & 40 & 576 & 546--689 \\
疾病治疗 & 肺炎有哪些治疗方法？ & 34 & 583 & 530--631 \\
就诊科室 & 肺炎应该挂哪个科室？ & 12 & 609 & 568--610 \\
常用药物 & 肺炎常用哪些药物？ & 40 & 824 & 747--852 \\
科室分布 & 疾病在各科室的分布情况如何？ & 30 & 570 & 537--588 \\
实体分布 & 医疗知识图谱中各类实体数量是多少？ & 12 & 553 & 540--580 \\
关系分布 & 知识图谱中不同关系类型的数量是多少？ & 16 & 539 & 535--558 \\
\bottomrule
\end{tabularx}
\end{table}

八个问题均由医疗语义层完成规划并返回结构化结果。
五类疾病属性查询生成文字回答和结构化结果表，
三类分布查询同时生成统计图。
连续请求的中位时间位于 539--824\,ms，
药物与症状查询各返回 40 行；
关系分布查询的观测中位数最低；表中数据用于描述总体响应，不拆分网络、规划和查询各阶段的耗时。

\section{任务三 NL2SQL 执行评测}

在线业务核验关注服务是否能完成一次端到端请求，
NL2SQL 评测则单独检查自然语言问题能否生成可执行查询并返回正确结果。
评测目录包含 133 道题：开发集 40 题，测试集 93 题；
题目覆盖单表筛选、聚合统计、分组统计、Top-K 排序、多条件查询、多表联查和复杂综合七类 SQL 结构。
构建脚本先在真实的 \texttt{task3\_analytics.db} 上执行 Gold SQL，
再保存自然语言问题和 Gold Result，
因此每道题都有可复核的查询结果和数据口径。

评测主指标为执行结果准确率：
系统生成的 SQL 在同一分析库执行后，
返回行按多重集合比较，行内重复记录保留，顺序差异不影响判定。
该指标回答的是“查询结果是否正确”，
不等价于 SQL 字符串相同，也不衡量知识库检索、来源引用、图表样式或 Nexent 完整对话质量。
复现入口为 \texttt{evaluation/task3/run\_benchmark.py}，
题目、数据库和评测输出分别保存在 \texttt{evaluation/task3/} 下。

\begin{table}[H]
\centering
\small
\caption{任务三 NL2SQL 执行准确率与评测口径}
\label{tab:task3-nl2sql-summary}
\begin{tabularx}{\textwidth}{L{3.0cm}C{1.7cm}C{1.7cm}Y}
\toprule
\textbf{运行对象} & \textbf{正确/总题数} & \textbf{执行准确率} & \textbf{口径说明} \\
\midrule
开发集初始方案 & 25/40 & 62.5\% & 开发集初始运行，用于识别模板覆盖和查询规划缺口 \\
开发集方案验证 & 38/40 & 95.0\% & 在开发集上验证语义模板与结构化计划的改动 \\
测试集设计基线 & 57/93 & 61.3\% & 测试集首次运行，未根据测试失败题型调整 \\
测试集回归复核 & 89/93 & 95.7\% & 已使用测试失败题型信息，属于回归复核，不是独立盲测 \\
\bottomrule
\end{tabularx}
\end{table}

表~\ref{tab:task3-nl2sql-summary} 说明，测试集设计基线低于赛题要求的 85\%，
回归复核达到 95.7\%。
回归复核中，聚合统计、分组统计、多条件查询、Top-K 排序和多表联查均达到 100\%；
复杂综合由 27.3\% 提升到 63.6\%，仍有 4 道题未通过。
这组结果支持“统一语义模板和受控查询计划能覆盖固定分析问题”的工程判断，
但不能把回归复核结果解释为对未见问题的泛化准确率。

\begin{table}[H]
\centering
\small
\caption{测试集按题型的 NL2SQL 执行准确率}
\label{tab:task3-nl2sql-by-type}
\begin{tabularx}{\textwidth}{L{2.8cm}C{1.1cm}C{1.8cm}C{1.8cm}Y}
\toprule
\textbf{题型} & \textbf{题数} & \textbf{设计基线} & \textbf{回归复核} & \textbf{结果解释} \\
\midrule
单表筛选 & 17 & 100.0\% & 100.0\% & 基本疾病属性查询覆盖稳定 \\
聚合统计 & 17 & 58.8\% & 100.0\% & 聚合字段和排序口径经模板固化 \\
分组统计 & 14 & 64.3\% & 100.0\% & 分组、排序和限制条件覆盖完整 \\
多条件查询 & 13 & 53.8\% & 100.0\% & 多条件过滤和子查询错误得到修正 \\
Top-K 排序 & 12 & 83.3\% & 100.0\% & 排序方向和限制条数得到统一 \\
多表联查 & 9 & 11.1\% & 100.0\% & 联接键和连接路径是初始方案的主要缺口 \\
复杂综合 & 11 & 27.3\% & 63.6\% & 多跳关联、比例和综合聚合仍是当前弱项 \\
\midrule
合计 & 93 & 61.3\% & 95.7\% & 回归复核结果含测试集失败信息 \\
\bottomrule
\end{tabularx}
\end{table}

\textbf{评测结论。}
任务三目前可以稳定处理评测集中的常见疾病属性和固定统计结构，
复杂综合问题仍需要人工检查查询计划和结果口径。
该评测证明的是分析库上的 SQL 执行一致性，
任务三完整链路还需要结合 Nexent 工具调用、来源解释和工作台展示进行独立验收。

\section{阶段耗时与性能控制}

三个阶段具有不同的性能特征。
任务一的时间主要由文件下载、DataMate 算子任务和 PDF 解析组成，
系统按格式独立组织子任务，文件数超过 50 时转入后台；
任务二按记录完成抽取与校验，
写入采用数据库事务，并且只在产生新三元组时刷新分析库；
任务三由统一分析服务读取本地分析数据库，
通过模板或结构化计划、只读连接、5 秒中止条件和 200 行上限控制单次负载。

\begin{table}[H]
\centering
\caption{三阶段耗时记录与控制方式}
\label{tab:performance-controls}
\begin{tabularx}{\textwidth}{L{2.3cm}L{5.3cm}Y}
\toprule
\textbf{阶段} & \textbf{主要耗时} & \textbf{控制方式} \\
\midrule
数据处理 &
文件传输、DataMate 算子、PDF 版面解析、结果汇聚 &
按格式拆分任务；大数据集后台执行；保存分组进度 \\
\addlinespace
知识构建 &
逐记录抽取、实体关系校验、图谱写入、分析表刷新 &
轮转取样；事务写入；有新增事实时刷新分析库 \\
\addlinespace
分析查询 &
语义规划、SQLite 查询、回答与图表数据组织 &
模板优先；最多 4 项查询；5 秒中止；最多 200 行 \\
\bottomrule
\end{tabularx}
\end{table}

当前八类在线查询的中位响应时间为 539--824\,ms，
本次观测的最大值为 1\,028\,ms；
该结果用于描述当前服务实例的响应情况。
任务一与任务二涉及外部平台、文档解析和知识抽取，
系统返回结构已分别记录分组耗时、逐记录耗时、入库数量和刷新耗时，
这些字段为任务一和任务二在相同数据规模、配置和部署条件下的后续比较提供统一口径。
